%% ============================================================
%% MODELO DE ATIVIDADE · COMUNICAÇÃO, MARCA E RELAÇÕES INSTITUCIONAIS
%%
%% Copie este arquivo para atividades/ (atividade-02.tex, por exemplo),
%% inclua-o no relatorio-aacc.tex com \input{atividades/atividade-02}
%% e escreva. Cada \preencher{…} é uma pendência: troque-o pelo seu
%% texto. Cada \figuraprovisoria também: troque-a pela figura de verdade.
%%
%% Uma atividade é UMA entrega sua, com começo, meio e fim. "Participei
%% do projeto X por dois anos" não é uma atividade: são várias. Divida.
%% ============================================================
\begin{atividade}{Campanha C de divulgação do projeto P}{
  tipo         = comunicacao,
  modalidade   = {},   % a chave da modalidade pedida: EXT, PES, COMP, VPC, EVT, PRD, ORG, ENS, GES, CUR, OUT
  alternativa  = {},   % a modalidade alternativa, se o Colegiado não aceitar a primeira
  horas        = {},   % as horas desta atividade, com base nos registros (Parte IV)
  inicio       = {},   % AAAA-MM
  fim          = {},   % AAAA-MM, ou: em andamento
  projeto      = {},
  papel        = {},   % Responsável, Corresponsável ou Colaborador(a)
  supervisor   = {},   % quem confirma: nome e cargo (o supervisor do projeto, o gerente)
  registros    = {},   % os códigos no SOMA, com ponto e vírgula: NBL-14; NRO-PRO-003-7
  competencias = {},   % as competências das DCN que ela desenvolveu (I a VIII): I, III, V
  disciplinas  = {},   % as disciplinas do seu curso ligadas a ela: ELE075 Eletrônica; ELE067 Circuitos
}

\begin{contexto}
\preencher{O objetivo de comunicação ou de relacionamento, e o público.}
\end{contexto}

\begin{contribuicao}
\preencher{O que você produziu ou negociou.}
\end{contribuicao}

\begin{desenvolvimento}
\fichatecnica{
  canais    = {},   % Canais e públicos
  pecas     = {},   % Peças e campanhas (códigos POST)
  parceiros = {},   % Parceiros e instituições — opcional
  metricas  = {},   % Métricas (alcance, engajamento, resultado)
}

\preencher{O planejamento, as peças, a aprovação, a publicação.}

\begin{figure}[htbp]
  \centering
  \figuraprovisoria{Uma peça produzida, ou o painel com as métricas.}
  % \includegraphics[width=0.9\linewidth]{figuras/peca.pdf}
  \caption{Peça da campanha C publicada no Instagram da equipe.}
  \fonte{Elaborado pelo autor.}
  \label{fig:peca}
\end{figure}
\end{desenvolvimento}

\begin{resultados}
\preencher{Os números: o alcance, o engajamento, as parcerias fechadas.}
\end{resultados}

\begin{evidencias}
% Uma linha por evidência: \evidencia{tipo}{identificador}{o que mostra}{onde conferir}
% \evidencia{Publicação}{POST-<n>}{A publicação aprovada e publicada}{SOMA › Studio}
\end{evidencias}

\begin{aprendizados}
\preencher{O que comunicar engenharia ensinou.}
\end{aprendizados}
\end{atividade}
